\documentclass[10pt,a4paper]{amsart}
\usepackage[margin=19mm]{geometry}
\usepackage{amsmath,amssymb,mathtools}
\usepackage[scr=rsfs,cal=euler]{mathalfa}
\usepackage{xfrac}
\usepackage[unicode,hidelinks,bookmarks=false]{hyperref}
\newcommand{\ie}{i.e.\ }
\newcommand{\cf}{cf.\ }
\newcommand{\resp}{respectively\ }
\newcommand{\Subsection}{\subsection}
\providecommand{\nobreakdash}{\nobreak\hbox{-}}
\setlength{\emergencystretch}{2em}
\multlinegap0pt
\usepackage{bm}
\usepackage{bbm}
\usepackage{dsfont}
\usepackage{xspace}
\usepackage{cancel}
\usepackage{mathtools}
\usepackage{amsthm}
\makeatletter
\@ifundefined{theorem}{\newtheorem{theorem}{Theorem}}{}
\makeatother
\title[Global solutions for systems of strongly invariant operators]{Global solutions for systems of strongly invariant operators on closed manifolds}
\author{Alexandre Kirilov, Wagner Augusto Almeida de Moraes, Pedro Meyer Tokoro}
\date{Source: arXiv:2502.18660v1 (CC BY 4.0)}
\begin{document}
\maketitle

\noindent\emph{Source and licence.} Global solutions for systems of strongly invariant operators on closed manifolds. Version arXiv:2502.18660v1 (CC BY 4.0), distributed under the Creative Commons Attribution 4.0 International licence, as recorded in the arXiv licence metadata for this exact version and shown on the abstract page https://arxiv.org/abs/2502.18660v1. Full rights notice: licenses/arxiv-2502.18660-CC-BY-4.0.txt.\par\smallskip

\noindent\textit{Editorial scope.} Counted unit: the Theorem whose source label is \texttt{equiv\_gs\_agh} in the article (source0.tex lines 344--346, source label \texttt{equiv\_gs\_agh}), delivered together with the article's own complete original proof of it (source0.tex lines 348--391, one \texttt{proof} environment of 44 lines). No mathematics is written, repaired or re-derived: statement and proof are the article's own text line for line. Required context retained uncounted: 6.1 (lines 262--264); agh\_est (lines 340--342). Every definition, display and label of those lines is retained unchanged. Omitted from this package and not redistributed: 1--261, 265--339, 343--343, 347--347, 392--855. Nothing inside a retained excerpt is omitted or altered: every retained body is delivered exactly as the article's lines are written, so no omission receipt of any kind accompanies this package. Citations: the retained excerpts carry no citation command at all, so no bibliography entry of the article is transcribed and no reference is redistributed. Reference adaptation: none was needed and none was made; every reference inside the delivered unit resolves inside it, so no unresolved reference or citation is produced. Printed slips kept literal: no printed word, symbol, hypothesis or number of the counted statement or of its proof was altered. Layout and class: the article's own class is \texttt{amsart} with packages including amssymb, indentfirst, enumerate, cite, amsfonts, amsmath, amsthm, mathrsfs, dsfont, hyperref; the wrapper is the lane's pinned \texttt{amsart} editorial wrapper at 10pt on A4, which restates the theorem-like environments the retained text needs and adds the article's own macro definitions that the retained text needs (none), with extra wrapper packages (bm, bbm, dsfont, xspace, cancel, mathtools). The delivered numbering is the unit's own. Assets: no retained span contains a figure environment, a graphics-inclusion command, a PDF-page inclusion, a table or a TikZ picture; no image file is copied into this package, the package declares no asset dependency and no image file is redistributed with it. Licence: version arXiv:2502.18660v1 of this article is distributed under the Creative Commons Attribution 4.0 International licence, as recorded in the arXiv licence metadata for this exact version and shown on the abstract page https://arxiv.org/abs/2502.18660v1. A full rights notice is delivered at licenses/arxiv-2502.18660-CC-BY-4.0.txt. Characters: every retained excerpt is pure ASCII, so the delivered engine, XeLaTeX with its default Latin Modern text font, reports no missing character.
		\begin{equation}\label{6.1}
			\|\sigma_P(k)\varphi\|\geq C(1+\lambda_k)^{t/\nu},\ \forall \varphi\in\ker\sigma_P(k)^\perp,\ \|\varphi\|_{L^2}=1,\ k\in\mathbb{N}_0.
		\end{equation}

	\begin{theorem}\label{agh_est}
		A strongly invariant operator $P$ is almost globally hypoelliptic if and only if there exist $s\in\mathbb{R}$ and $C>0$ such that condition (\ref{6.1}) holds.
	\end{theorem}

	\begin{theorem}\label{equiv_gs_agh}
		A strongly invariant operator $P$ is globally solvable if and only if it is almost globally hypoelliptic.
	\end{theorem}

	\begin{proof}
		\noindent($\Leftarrow$) Suppose that $P$ is almost globally hypoelliptic.
		Let $\{u_\ell\}_{\ell\in\mathbb{N}}$ be a sequence in $C^\infty(M)$ and let $f\in C^\infty(M)$ be such that $Pu_\ell\to f$ in $C^\infty(M)$. Then,
		\begin{equation*}
			\lim\limits_{\ell\to\infty}\sigma_P(k)\widehat{u}_\ell(k) = \widehat{f}(k), \text{ for all } k\in\mathbb{N}_0.
		\end{equation*}
		
		Since each $\sigma_P(k):E_{\lambda_k}\to E_{\lambda_k}$ has a closed image (as $E_{\lambda_k}$ is finite-dimensional), there exists $\widehat{u}(k)$ such that $\sigma_P(k)\widehat{u}(k)=\widehat{f}(k)$. Furthermore, we can choose $\widehat{u}(k)$ such that $\widehat{u}(k) \in \ker\sigma_P(k)^\perp$.  
		
		By Theorem \ref{agh_est}, there exist constants $C>0$ and $t\in\mathbb{R}$ such that
		\begin{equation*}
			\|\widehat{u}(k)\|\leq C^{-1}(1+\lambda_k)^{-t/\nu}\|\widehat{f}(k)\|.
		\end{equation*}

		Since $f\in C^\infty(M)$, this bound implies that $u\in C^\infty(M)$. Moreover, we have $Pu=f$, which shows that $f\in P(C^\infty(M))$. Hence, $P$ is globally solvable.  
		
		\medskip\noindent
		($\Rightarrow$) Suppose that $P$ is globally solvable, i.e., $P(C^\infty(M))$ is closed.
		We first show that the space $\mathcal{F}=\{u\in C^\infty(M)\,:\, \widehat{u}(k)\in\ker\sigma_P(k)^\perp,\ \forall k\in\mathbb{N}_0\}$	is closed in $C^\infty(M)$. 
		
		Indeed, if $\{u_\ell\}_{\ell\in\mathbb{N}}$ is a sequence such that $u_\ell\to u$ in $C^\infty(M)$, then $\widehat{u}_\ell(k)\to\widehat{u}(k)$ for all $k\in\mathbb{N}_0$. Since each $\ker\sigma_P(k)^\perp$ is closed in $E_{\lambda_k}$, it follows that $u\in\mathcal{F}$.  
		
		Since $P:\mathscr{D}'(M)\to\mathscr{D}'(M)$ is continuous, the Closed Graph Theorem implies that the restriction $P:C^\infty(M)\to C^\infty(M)$ is also continuous. This, in turn, implies that the induced operator $P:\mathcal{F}\to P(C^\infty(M))$ is a continuous linear bijection between FS spaces. By the Open Mapping Theorem, its inverse  $P^{-1}: P(C^\infty(M))\to\mathcal{F}$ is also continuous.  
		
		By the definition of the projective limit topology, this continuity means that for each $t>0$, there exist $s>0$ and $C>0$ such that
		\begin{equation*}
			\|P^{-1}v\|_{H^t}\leq C\|v\|_{H^s},\quad \forall v\in P(C^\infty(M)).
		\end{equation*}
		Equivalently, for all $u\in\mathcal{F}$,
		\begin{equation*}
			\|u\|_{H^t}\leq C\|Pu\|_{H^s}.
		\end{equation*}
		
		In particular, for each $k\in\mathbb{N}_0$ and $\varphi\in\ker\sigma_P(k)^\perp$, we have $\varphi\in\mathcal{F}$, so
		\begin{equation*}
			(1+\lambda_k)^t\|\varphi\|=\|\varphi\|_{H^t}\leq C\|\sigma_P(k)\varphi\|_{H^s}=C(1+\lambda_k)^s\|\sigma_P(k)\varphi\|.
		\end{equation*}
		This implies that
		\begin{equation*}
			\|\sigma_P(k)\varphi\|\geq C^{-1}(1+\lambda_k)^{t-s}\|\varphi\|,\quad \forall \varphi\in\ker\sigma_P(k)^\perp, \quad \forall k\in\mathbb{N}_0.
		\end{equation*}
		
		By Theorem \ref{agh_est}, this establishes that $P$ is almost globally hypoelliptic.
	\end{proof}

\end{document}
